% !TeX root = ../main.tex

% ==================================================
% 1.4. import, module, package, sys.path
% ==================================================


\subsection{import, module, package, sys.path}

우리는 당연하게 아래와 같은 것을 사용한다.

\begin{dhkcode}
import argparse
from pathlib import Path
import pandas as pd
import torch
\end{dhkcode}

처음에는 그냥:

\begin{dhksmry}[centered]
라이브러리를 가져온다.
\end{dhksmry}

정도로 배운다.

그런데 이제는 그 설명을 버리고 좀 더 정확하게 볼 때이다.

예를 들어:

\begin{dhkcode}
import pandas
\end{dhkcode}

라고 하면 Python은 대략적으로

\begin{dhksmry}[centered]
"pandas라는 이름의 module/package가 필요하네."
\da
어디에 있는지 찾아야 함
\da
import 탐색 경로
\da
sys.path
\da
site-packages/
\da
pandas를 발견
\da
module을 load/execute
\da
현재 namespace에서 pandas 라는 이름으로 사용
\end{dhksmry}

하는 과정이 있다.

여기서 우리가 앞에서 배운 \textbf{파일 경로}와 드디어 다시 만난다.

\subsubsection{sys.path를 먼저 보자}

Python에서:

\begin{dhkcode}
import sys

print(sys.path)
\end{dhkcode}

를 해보면 대략 이런 리스트가 나온다.

\begin{dhkzsh}
[
'/Users/user/Desktop/my_project/src',
'/Users/user/anaconda3/envs/transformer/lib/python310.zip',
'/Users/user/anaconda3/envs/transformer/lib/python3.10',
'/Users/user/anaconda3/envs/transformer/lib/python3.10/site-packages'
]
\end{dhkzsh}

중요한 것은 정확한 문자열을 외우는 게 아니다.

\textbf{sys.path 자체가 리스트다}.

\begin{dhkdir}
sys.path
>어떤 디렉터리
>어떤 디렉터리
>어떤 디렉터리
>site-packages
\end{dhkdir}

Python이:

\begin{dhkcode}
import pandas
\end{dhkcode}

를 만나면 이런 \textbf{탐색 위치들을 이용해서 pandas를 찾는다}.

그래서 아주 중요한 구분이 생긴다.

\begin{dhksmry}[centered]
Path.cwd()
\da
"현재 어디에서 작업하고 있는가?"

sys.path()
\da
"import할 module/package를 어디에서 찾을 것인가?"
\end{dhksmry}

둘 다 경로가 등장하니까 처음에는 굉장히 헛갈리지만 \textbf{질문 자체가 다르다}.

\subsubsection{그리고 \zsh{site-packages}가 등장한다}

Conda 환경에서:

\begin{dhkzsh}
pip install pandas
\end{dhkzsh}

했다고 해보자.

패키지는 보통 현재 Python 환경의 \zsh{site-packages}쪽에 설치된다.

개념적으로:

\begin{dhkdir}
transformer/
>lib/
>>python 3.10/
>>>site-packages/
>>>>pandas/
>>>>numpy/
>>>>transformers/
>>>>...
\end{dhkdir}

그리고 그 위치가 Python의 import탐색 대상에 있으니까:

\begin{dhkcode}
import pandas
\end{dhkcode}

가 가능하다.

전에 확인했던:

\begin{dhkcode}
sys.executable
\end{dhkcode}

과 연결하면 상당히 중요한 그림이 완성된다.

\begin{dhkdir}
conda transformer 환경
> Python interpreter
>> sys.executable
> site-packages
>> torch
>> pandas
>> transformers
\end{dhkdir}

그래서 VS Code/Jupyter에서 흔히 발생하는:

\begin{dhksmry}[centered]
분명히 \zsh{pip install} 했는데 \zsh{ModuleNotFoundError}가 나는데?
\end{dhksmry}

라는 문제도 결국 상당수가:

\begin{dhksmry}[centered]
설치한 Python 환경 $\neq$ 현재 코드를 실행하는 Python 환경
\end{dhksmry}

문제다.

\begin{dhktable}
|명령어|의미|
|:---|:---|
|Path.cwd()|내가 지금 어디서 작업하는가?|
|sys.executable|어떤 Python이 나를 실행하고 있는가?|
|sys.path|그 Python이 import 대상을 어디서 찾는가?|
|site-packages|설치한 외부 package들이 주로 존재하는 곳|
\end{dhktable}

이제 직접만든 \zsh{util.py}를 Python이 어떻게 찾아서 import하는지 볼 것이다.

여기서부터:

\begin{dhksmry}[centered]
module \ra package \ra sys.path \ra ModuleNotFoundError
\end{dhksmry}

가 한 번에 연결된다.

\subsubsection{먼저 module이 뭔가?}

아주 단순하게 시작하면 \textbf{Python 파일 하나가 하나의 module이 될 수 있다}.

프로젝트가:

\begin{dhkdir}
my_project/
>train.py
>utils.py
\end{dhkdir}

이렇게 있다고 하자.

\zsh{utils.py}:

\begin{dhkcode}
def add(a, b):
    return a + b
\end{dhkcode}

\zsh{train.py}:

\begin{dhkcode}
import utils

results = utils.add(10,20)
print(result)
\end{dhkcode}

여기서:

\begin{dhkcode}
import utils
\end{dhkcode}

는 대충:

\begin{dhksmry}[centered]
\zsh{utils}라는 module을 찾아서 사용할 수 있게 해줘.
\end{dhksmry}

라는 뜻이다.

Python은 \zsh{utils.py}를 발견하면 그 파일을 module 객체로 만들어 다룬다.

\begin{dhkbranch}
node main :: utils.py
arrow main :: import
node main :: module 객체
arrow main
node main :: 이름 utils에 binding
\end{dhkbranch}

\begin{dhkdir}
utils
>add
\end{dhkdir}

그래서:

\begin{dhkcode}
utils.add(10, 20)
\end{dhkcode}

처럼 접근할 수 있다.

이것도:

\begin{dhksmry}[centered]
name \ra object binding
\end{dhksmry}

과 똑같다.

\begin{dhkdir}
import utils

이름        객체

util -> @module object
>add
\end{dhkdir}

\subsubsection{\zsh{from utils import add}는 뭐가 다를까?}

이번에는:

\begin{dhkcode}
from utils import add

result = add(10, 20)
\end{dhkcode}

라고 해보자.

개념적으로는 최종 namespace가 다르다.

\textbf{import utils}

\begin{dhkdir}
현재 namespace

utils -> @module
>add
\end{dhkdir}

따라서:

\begin{dhkcode}
utils.add()
\end{dhkcode}

\textbf{from utils import add}

\begin{dhksmry}
현재 namespace

add \ra function object
\end{dhksmry}

따라서:

\begin{dhkcode}
add()
\end{dhkcode}

이렇게 사용한다.

여기서도 완전히 새로운 원리가 나온 게 아니라 \textbf{객체와 이름의 binding}이 나온 것이다.

\subsubsection{그런데 Python은 \zsh{utils.py}를 어떻게 찾았나?}

이것이 이번 파트의 핵심이다.

\begin{dhkcode}
import utils
\end{dhkcode}

를 실행했다고 해서 Python이 컴퓨터 전체 SSD를 뒤지는 게 아니다.

Python에는 \textbf{import를 찾을 위치들}이 있고, 그게 앞에서 본:

\begin{dhkcode}
sys.path
\end{dhkcode}

이다.

개념적으로:

\begin{dhksmry}[centered]
import utils
\da
"utils 어디 있지?"
\da
import system
\da
sys.path에 있는 위치들을 탐색
\da
utils를 발견
\da
load
\da
module object 생성/사용
\da
utils라는 이름에 binding
\end{dhksmry}

여기까지가 첫 번째 큰 그림이다.

\subsubsection{이제 폴더가 하나 생기면 문제가 재미있어진다}

프로젝트를:

\begin{dhkdir}
my_project/
>train.py
>src/
>>utils.py
\end{dhkdir}

이렇게 바꿔보자.

이제 \zsh{train.py}에서:

\begin{dhkcode}
import utils
\end{dhkcode}

하면 \zsh{src/utils.py}를 찾을까?

\textbf{아니다}.

\zsh{src/}가 존재한다고 해서 

Python이 자동으로 모든 하위 디렉터리를 import탐색 대상으로 삼는 게 아니다.

이게 굉장히 중요하다.

Python이 \zsh{my_project}를 탐색한다고 해서:

\begin{dhksmry}[centered]
my_project
\da
모든 하위 폴더를 재귀적으로 검색
\da
src/utils.py 발견
\end{dhksmry}

하는 방식이 \textbf{아니다}.

그래서 import를 이해할 때:

\begin{dhksmry}[centered]
파일이 프로젝트 어딘가에 존재한다.
\end{dhksmry}

와:

\begin{dhksmry}[centered]
Python이 그 파일을 import대상으로 찾을 수 있다.
\end{dhksmry}

는 전혀 다른 말이다.

\subsubsection{그래서 package라는 개념이 필요해진다}

우리가 원하는 것은:

\begin{dhkcode}
from src import utils
\end{dhkcode}

또는:

\begin{dhkcode}
from src.utils import add
\end{dhkcode}

같은 구조라면 \zsh{src}를 \textbf{package}구조로 다루게 된다.

전통적으로 아주 흔하게 보는 구조가:

\begin{dhkdir}
my_project/
>train.py
>src/
>>__init__.py
>>utils.py
\end{dhkdir}

이다.

그러면 개념적으로:

\begin{dhkdir}
src
>__init__.py
>utils.py
>>add()
\end{dhkdir}

라는 \zsh{package/module}계층이 만들어진다.

그래서:

\begin{dhkcode}
from src.utils import add
\end{dhkcode}

를 보면 이제 이렇게 읽으면 된다.

\begin{dhkdiagram}[in=8]
node full (1,2) :: src.utils.add
node n1 (2,1)
node n2 (2,2)
node n3 (2,3)
node src (3,1) :: src
node util (3,2) :: utils
node add (3,3) :: add
node f (4,4) :: 함수/객체
node m (5,3) :: module
node p (6,2) :: package

edge full -- n2
edge n1 -- n2
edge n2 -- n3
edge src -- n1
edge n3 -- add
edge n2 -- util
edge add -- f :: VH
edge util -- m :: VH
edge src -- p :: VH
\end{dhkdiagram}

사용:

\begin{dhkcode}
from src.utils import add

print(add(10, 20))
\end{dhkcode}

\subsubsection{\zsh{__init__.py}는 뭔데?}

이 파일을 처음 보면 이상하다.

\begin{dhksmry}[centered]
__init__.py
\end{dhksmry}

클래스에서 배운:

\begin{dhkcode}
def __init__(self):
\end{dhkcode}

와 이름이 같아서 더 헛갈릴 수 있지만

지금은:

\begin{dhksmry}[centered]
package의 __init__.py
\end{dhksmry}

이다.

역사적으로 Python에서는:

\begin{dhkdir}
src/
>__init__.py
>utils.py
\end{dhkdir}

처럼:

\begin{dhksmry}[centered]
__init__.py
\end{dhksmry}

가 있으면 이 디렉터리를 Python package로 취급하는 전통적인 구조를 사용했다.

현대 Python에는 이 것이 없어도 가능한 \textbf{namespace package}도 있어서,

\begin{dhksmry}[centered]
package라면 무조건 __init__.py가 있어야 한다.
\end{dhksmry}

라고 외우면 정확하지 않다.

하지만 일반적인 Python 프로젝트 구조를 배우는 지금 단계에서는:

\begin{dhksmry}[centered]
디렉터리를 일반적인 Python package로 구성
\da
__init__.py를 둔다.
\end{dhksmry}

정도로 먼저 이해하면 충분하다.

\subsubsection{module과 package를 이렇게 구분하면 된다}

지금 단계에서는:

\begin{dhksmry}[centered]
module
\da
Python 코드의 한 단위
대표적으로 .py 파일

utils.py
train.py
model.py
\end{dhksmry}

반면:

\begin{dhksmry}[centered]
package
\da
여러 module 등을 계층적으로 조직하는 구조
\end{dhksmry}

\begin{dhkdir}
src/
>__init__.py
>utils.py
>model.py
>dataset.py
\end{dhkdir}

그래서 ML 프로젝트도 점점 커지면:

\begin{dhkdir}
my_project/
>train.py
>src/
>>__init__.py
>>data.py
>>model.py
>>utils.py
\end{dhkdir}

같아지고,

\begin{dhkcode}
from src.data import load_data
from src.model import MyModel
from src.utils import save_model
\end{dhkcode}

같은 코드가 나온다.

\subsubsection{여기서 \zsh{ModuleNotFoundError}의 의미가 보이기 시작한다}

예를 들어:

\begin{dhkcode}
from src.model import MyModel
\end{dhkcode}

했는데:

\begin{dhkzsh}
ModuleNotFoundError
\end{dhkzsh}

가 발생했다면 무조건:

\begin{dhksmry}[centered]
\zsh{model.py}파일이 없다.
\end{dhksmry}

라는 뜻이 아니다.

더 정확하게는:

\begin{dhknote}
\begin{dhksmry}[centered]
현재 import system이 주어진 이름을 해석해서

해당 module/package를 찾지 못했다.
\end{dhksmry}
\end{dhknote}

이다.

따라서 문제를 볼 때 이제 세 가지를 분리해야 한다.

\begin{dhksmry}[centered]
실제 파일은 어디 있는가?
\da
Path / filesystem 문제



현재 어떤 Python으로 실행 중인가?
\da
sys.executable / Conda 문제



그 Python은 import 대상을 어디서 찾는가?
\da
sys.path / package 구조 문제
\end{dhksmry}

이 구분이 중요하다.

\subsubsection{그리고 여기서 \zsh{sys.path.append()}가 왜 인터넷에 자주 나오는지도 보인다}

검색하다보면 이런 코드를 엄청 많이 보게 된다.

\begin{dhkcode}
import sys

sys.path.append("...")
\end{dhkcode}

이게 하는 일은 말 그대로:

\begin{dhksmry}
[
    A,
    B,
    C
]

append(D)

[
    A,
    B,
    C,
    D
]
\end{dhksmry}

즉:

\begin{dhksmry}[centered]
Python아, import할 때 이 디렉터리도 찾아봐.
\end{dhksmry}

라고 import 탐색 위치를 런타임에 추가하는 것이다.

그래서 import가 안 될 때:

\begin{dhkcode}
sys.path.append("...")
\end{dhkcode}

같은 것을 넣으면 갑자기 되는 경우가 있다.

그런데 이건 때때로 \textbf{프로젝트 구조나 실행 방법의 문제를 임시로 덮어버리는 것}일 수 있다.

그래서 지금 단계에서는:

\begin{dhksmry}[centered]
ModuleNotFoundError 가 발생했으니 sys.path.append부터 넣자
\end{dhksmry}

가 아니라

\begin{dhksmry}[centered]
ModuleNotFoundError
\da
파일은 어디 있지?
\da
어떤 Python으로 실행 중이지?
\da
어디서 어떤 방식으로 실행했지?
\da
package 구조는 어떻게 되어 있지?
\da
sys.path는 어떻게 구성되어 있지?
\end{dhksmry}

라고 생각하는 습관을 들이는 것이 좋다.

\subsubsection{여기까지의 핵심}

새로운 내용을 기존 개념과 합치면 이렇게 된다.

\begin{dhkdir}[in=6]
File system
> utils.py → module
> src/
>>__init__.py → package 구성
>>model.py → module
\end{dhkdir}

\begin{dhksmry}[centered]
import src.model
\da
Python import system
\da
sys.path
\da
src를 찾음
\da
...
\end{dhksmry}

이것이 \textbf{Filesystem + Path + name/object/binding + 실행환경}이 하나로 합쳐지는 지점이다.

다음 단계에서는 가장 많이 헛갈리는

\begin{dhkzsh}
python src/train.py

python -m src.train
\end{dhkzsh}

위 두 가지가 서로 다르고, 

왜 하나에서는 import가 되지만 다른 하나에서는 깨질 수 있는지

를 실제 프로젝트 구조를 놓고 볼 것이다.

\subsubsection{같은 \zsh{train.py}인데 실행 방법이 두 가지다}

프로젝트를 이렇게 두자:

\begin{dhkdir}
my_project/
>src/
>>__init__.py
>>train.py
>>utils.py
>data/
>>train.csv
\end{dhkdir}

그리고 \zsh{train.py}에서 같은 \zsh{package}의 \zsh{utils.py}를 사용하고 싶다고 하자.

\begin{dhkcode}
from src.utils import add
\end{dhkcode}

터미널의 CWD는:

\begin{dhkzsh}
my_project/
\end{dhkzsh}

라고 하자.

여기서 두 가지 실행법이 있다.

\begin{dhkzsh}
python src/train.py
\end{dhkzsh}

그리고:

\begin{dhkzsh}
python -m src.train
\end{dhkzsh}

둘 다 \zsh{train.py}의 코드를 실행하는 것 처럼 보이는가?

그렇지만 Python에게 전달하는 \textbf{의미가 다르다}.

\subsubsection{\zsh{python src/train.py}}

이건 우리가 지금까지 많이 썼던 방식이다.

\begin{dhkzsh}
python src/train.py
\end{dhkzsh}

Python에게 말하는 건:

\begin{dhksmry}[centered]
\zsh{src/train.py}라는 \textbf{파일을 직접 실행해}.
\end{dhksmry}

라는 의미다.

즉 출발점이 \textbf{파일 경로}이다.

\begin{dhkdir}
python src/train.py
>>>>파일 경로
\end{dhkdir}

이때 \zsh{train.py}는 실행할 최상위 스크립트로 취급된다.

그리고 중요한 변화가 하나 생긴다.

실행되는 스크립트가 있는 디렉터리인:

\begin{dhkzsh}
my_project/src/
\end{dhkzsh}

가 import탐색에 중요한 위치로 들어오게 된다.

그래서 상황에 따라 package import가 예상과 다르게 동작할 수 있다.

\subsubsection{\zsh{python -m src.train}}

이번엔 완전히 관점이 다르다.

\begin{dhkzsh}
python -m src.train
\end{dhkzsh}

여기에는:

\begin{dhkzsh}
src/train.py
\end{dhkzsh}

라는 파일 경로를 주지 않았다.

대신:

\begin{dhkzsh}
src.train
\end{dhkzsh}

이라는 \textbf{module 이름}을 주었다.

\zsh{-m}은 간단히 말하면:

\begin{dhksmry}[centered]
이 이름의 module을 찾아서 실행해.
\end{dhksmry}

라는 의미다.

따라서 Python의 사고 과정은:

\begin{dhksmry}[centered]
python -m src.train

현재 import 환경에서

\da

src라는 package를 찾고

\da

그 안의 train module을 찾고

\da

그 module을 실행
\end{dhksmry}

이렇게 된다.

즉:

\begin{dhksmry}[centered]
python src/train.py

\da

파일 경로를 직접 실행



python -m src.train

\da

module/package 이름으로 찾아서 실행
\end{dhksmry}

이 차이가 핵심이다.

\subsubsection{왜 /를 안 쓰고 .을 쓰는가?}

이제 이유가 보일 것이다.

\begin{dhkzsh}
python src/train.py
\end{dhkzsh}

에서는 \textbf{Filesystem Path}니까:

\begin{dhkzsh}
src/train.py
\end{dhkzsh}

라고 쓴다.

반면:

\begin{dhkzsh}
python -m src.train
\end{dhkzsh}

에서는 \textbf{Python module}경로니까:

\begin{dhkzsh}
src.train
\end{dhkzsh}

이라고 쓴다.

따라서 둘을 구별해야 한다.

\begin{dhksmry}
    Filesystem
    ───────────────────────
    src/train.py
    src/utils.py



    Python module namespace
    ───────────────────────
    src.train
    src.utils
\end{dhksmry}

이것은 상당히 중요한 구분이다.

\subsubsection{이제 package 내부 import가 나온다}

\zsh{train.py}와 \zsh{utils.py}가 같은 \zsh{src} package 안에 있다.

\begin{dhkdir}
src/
>__init__.py
>train.py
>utils.py
\end{dhkdir}

\zsh{train.py}에서 이렇게 쓸 수도 있다.

\begin{dhkcode}
from src.utils import add
\end{dhkcode}

이걸 \textbf{absolut import}라고 생각하면 된다.

전체 package 기준으로:

\begin{dhkdir}
src
>utils
>>add
\end{dhkdir}

를 지정했기 때문이다.

그런데 이렇게도 쓸 수 있다.

\begin{dhkcode}
from .utils import add
\end{dhkcode}

여기서 . 은 Filesystem의:

\begin{dhkzsh}
./
\end{dhkzsh}

와 똑같은 문자열 문법이라고 생각하면 \textbf{안 된다}.

여기서는 \textbf{현재 package}라는 의미다.

즉 \zsh{train.py}가 \zsh{src} package 안에 있으므로:

\begin{dhkcode}
from .utils import add
\end{dhkcode}

는 개념적으로:

\begin{dhksmry}[centered]
현재 package

\da

src

\da

utils

\da

add
\end{dhksmry}

라는 뜻이다.

이를 \textbf{relative import}라고 한다.

\subsubsection{그런데 파일을 직접 실행하면 문제가 생길 수 있다}

\zsh{train.py}:

\begin{dhkcode}
from .tuils import add

print(add(10, 20))
\end{dhkcode}

라고 해보자.

그리고:

\begin{dhkzsh}
ptyhon src/trian.py
\end{dhkzsh}

로 실행하면 \zsh{relative import}관련 오류를 만날 수 있다.

왜냐하면 우리는 Python에게:

\begin{dhksmry}[centered]
\zsh{src} package 안의 \zsh{train} module을 실행해.
\end{dhksmry}

라고 한 것이 아니다.

그냥:

\begin{dhksmry}[centered]
이 \zsh{train.py}를 직접 실행해.
\end{dhksmry}


라고 했기 때문이다.

즉 Python입장에서는 \zsh{train.py}가 package 안에서

어떤 위치를 갖는 module인지에 대한 문맥이 깨질 수 있다.

반대로:

\begin{dhkzsh}
python -m src.train
\end{dhkzsh}

이라고 하면 처음부터

\begin{dhksmry}[centered]
src package

\da

train module
\end{dhksmry}

이라는 package 문맥으로 찾았다.

따라서:

\begin{dhkcode}
from .utils import add
\end{dhkcode}

의 . 도:

\begin{dhksmry}[centered]
현재 \zsh{package} = \zsh{src}
\end{dhksmry}

라고 해석할 수 있다.

이게 \zsh{-m}이 중요한 이유 중 하나이다.

\subsubsection{\zsh{__name__}까지 여기서 연결된다}

Python 파일에서 많이 보는:

\begin{dhkcode}
if __name__ == "__main__":
    main()
\end{dhkcode}

도 사실 이 주제이다.

Python이 어떤 module을 \textbf{실행의 시작점}으로 실행하면

그 모듈의:

\begin{dhkzsh}
__name__
\end{dhkzsh}

을:

\begin{dhkzsh}
"__main__"
\end{dhkzsh}

으로 설정한다.

예를 들어:

\begin{dhkcode}
def main():
    print("training.")

if __name__ == "__main__":
    main()
\end{dhkcode}

직접 실행하면:

\begin{dhkzsh}
python train.py
\end{dhkzsh}

\begin{dhksmry}[centered]
__name__=="__main__" 이므로 main()이 실행된다.
\end{dhksmry}

반면 다른 코드에서:

\begin{dhkcode}
import train
\end{dhkcode}

하면 보통:

\begin{dhkzsh}
__name__=="train"
\end{dhkzsh}

이므로 저 조건은 \zsh{False}가 된다.

그래서 \textbf{import 했을 때 학습이 갑자기 시작되는 것을 막고, 실행 진입점만 분리}할 수 있다.

\subsubsection{이게 ML 코드에서 왜 중요할까?}

프로젝트가 커지면 보통:

\begin{dhkdir}
my_project/
>src/
>>__init__.py
>>train.py
>>model.py
>>dataset.py
>>utils.py
>data/
>models/
\end{dhkdir}

이런 식이 된다.

\zsh{train.py} 혼자 모든 걸 하는 게 아니라:

\begin{dhkcode}
from src.model import MyModel
from src.dataset import load_dataset
from src.utils import save_model
\end{dhkcode}

처럼 여러 module이 연결된다.

그러면 단순히:

\begin{dhksmry}[centered]
파일 A 실행
\end{dhksmry}

만 생각하는 것보다:

\begin{dhkdir}
my_project
| src package
> train module
> model module
> dataset module
> utils module
\end{dhkdir}

이라는 package 구조로 생각하는 게 중요해진다.

그래서 프로젝트 루트에서:

\begin{dhkzsh}
python -m src.train
\end{dhkzsh}

처럼 실행하는 방식이 자연스러워 지는 경우가 있다.

\begin{dhkdir}
Python 프로젝트
> file system
>> src/train.py - Path
>> src/utils.py - Path
> import system
>> src.train - sys.path
>> src.utils - sys.path
>>> python src/train.py - 파일로 실행
>>> python -m src.train - module로 찾아 실행
\end{dhkdir}

여기까지 이해되면 \textbf{import의 뼈대}는 거의 잡힌 것이다.

다음 내용은 여기서 자연스럽게:

\begin{dhksmry}
    __name__

    __main__

    import 할 때 Python 파일의 코드가 실제로 실행된다는 사실

    module caching인 sys.modules
\end{dhksmry}

로 넘어가는 것이 좋다.

이것을 배우면:
\begin{dhknote}
\begin{dhksmry}[centered]
왜 import는 단순히 함수/클래스 정의를 복사해오는 게 아닌가?
\end{dhksmry}
\end{dhknote}

까지 이해할 수 있다.

\subsubsection{import는 정의를 복사해오는 게 아니다}

예를 들어 \zsh{utils.py}가 있다고 하자.

\begin{dhkcode}
print("utils.py 실행됨")

def add(a, b):
    return a + b
\end{dhkcode}

그리고 \zsh{train.py}:

\begin{dhkcode}
import utils

print(utils.add(10,20))
\end{dhkcode}

실행하면:

\begin{dhkzsh}
utils.py 실행됨
30
\end{dhkzsh}

이게 중요한 이유는,

\begin{dhkcode}
import utils
\end{dhkcode}

가 단순히 \zsh{add()}라는 함수 정의만 가져오는 게 아니라는 것이다.

Python은 처음 import 할 때 개념적으로:

\begin{dhksmry}[centered]
import utils

\da

utils라는 모듈을 찾음

\da

utils.py의 코드를 실행

\da

그 실행으로 module namespace 구성

\da

module 객체를 만듦

\da

utils라는 이름에 그 module 객체를 binding
\end{dhksmry}

한다고 보면 된다.

\subsubsection{왜 파일을 실행해야 할까?}

\zsh{utils.py}가:

\begin{dhkcode}
x = 10

def add(a, b):
    return a + b

class Calculator:
    pass
\end{dhkcode}

라고 해보자.

Python이 이 파일을 실행하면서 namespace가 만들어진다고 생각하면 된다.

\begin{dhksmry}
    utils module의 namespace

    x \ra 10

    add \ra function object

    Calculator \ra class object
\end{dhksmry}

그리고:

\begin{dhkcode}
import utils
\end{dhkcode}

하는 파이썬 파일에서는:

\begin{dhkdir}
현재 namespace

utils -> @module object
>x
>add
>Calculator
\end{dhkdir}


가 되는 것이다.

그래서:

\begin{dhkcode}
utils.x
utils.add()
utils.Calculator()
\end{dhkcode}

처럼 접근할 수 있었던 것이다.

이건 예전에 배운 \textbf{name $\rightarrow$ object binding}이 module 수준으로 확장된 것이다.

\subsubsection{그러면 import 할 때 위험한 코드도 실행된다}

예를 들어 \zsh{train.py}를 이렇게 만들면:

\begin{dhkcode}
def train():
    print("모델 학습 시작")

train()
\end{dhkcode}

그리고 다른 파일에서:

\begin{dhkcode}
import train
\end{dhkcode}

만 해도:

\begin{dhkzsh}
모델 학습 시작
\end{dhkzsh}

이 출력된다.

왜냐하면 Python은 \zsh{train.py}를 import 하기 위해

\textbf{그 모듈의 코드를 실행했기 때문}이다.

바로 여기서 아까 배운 것이 다시 등장한다.

\begin{dhkcode}
def main():
    print("모델 학습 시작")

if __name__ == "__main__":
    main()
\end{dhkcode}

이제 다른 파일에서:

\begin{dhkcode}
import train
\end{dhkcode}

하면 \zsh{train.py} 자체는 import 과정에서 실행된다.

그래서:

\begin{dhkcode}
def main():
\end{dhkcode}

도 실행되어 \textbf{함수 객체가 만들어지고 main이라는 이름에 binding}된다.

하지만:

\begin{dhkcode}
if __name__ == "__main__":
\end{dhkcode}

에서는 import 된 모듈이므로:

\begin{dhksmry}[centered]
__name__ == "train"

"train" == "__main__"
\da
False
\end{dhksmry}

따라서:

\begin{dhkcode}
main()
\end{dhkcode}

이라는 \textbf{함수 호출만 실행되지 않은 것}이다.

이 차이가 중요하다.

\begin{dhksmry}[centered]
import train

\da

train.py 자체가 실행되지 않는다 (X)

train.py는 실행된다 (O)

하지만 __main__ 조건 내부의 main() 호출이 실행되지 않는다. (O)
\end{dhksmry}

\subsubsection{그렇다면 매번 import할 때마다 파일을 다시 실행할까?}

여기서 \zsh{sys.modules}가 등장한다.

\zsh{utils.py}:

\begin{dhkcode}
print("UTILS EXECUTED")
\end{dhkcode}

\zsh{train.py}:

\begin{dhkcode}
import utils
import utils
import utils
\end{dhkcode}

얼핏 생각하면:

\begin{dhkzsh}
UTILS EXECUTED
UTILS EXECUTED
UTILS EXECUTED
\end{dhkzsh}

일 것 같지만 일반적으로 한 번만 출력된다.

왜냐하면 Python은 이미 import한 module을 \textbf{캐시, Cache}해두기 때문이다.

그 캐시가:

\begin{dhksmry}[centered]
sys.modules
\end{dhksmry}

이다.

\subsubsection{\zsh{sys.modules}는 dictionary다}

\begin{dhkcode}
import sys

print(type(sys.modules))
\end{dhkcode}

개념적으로:

\begin{dhkzsh}
dict
\end{dhkzsh}

이고 대충 이런 구조라 생각하면 된다.

\begin{dhkzsh}
sys.modules
{
    "sys": <module 'sys'>,
    "os": <module 'os'>,
    "pathlib": <module 'pathlib'>,
    "utils": <module 'utils'>,
    ...
}
\end{dhkzsh}

즉:

\begin{dhksmry}[centered]
module 이름

\da

module 객체
\end{dhksmry}

를 저장하는 dictionary이다.

\subsubsection{두 번재 \zsh{import utils}에서는 어떻게 되는가?}

첫 번재:

\begin{dhkcode}
import utils
\end{dhkcode}

에서는 대략:

\begin{dhksmry}[centered]
import utils

\da

sys.modules에 utils 있나?

\da

없음

\da

모듈 탐색

\da

utils.py 발견

\da

코드 실행

\da

module 객체 생성/초기화

\da

sys.modules["utils"]에 등록
\end{dhksmry}

이제:

\begin{dhksmry}
    sys.modules

    "utils" \ra utils module object
\end{dhksmry}

가 되었다고 생각하자.

두 번재:

\begin{dhkcode}
import utils
\end{dhkcode}

에서는:

\begin{dhksmry}[centered]
import utils

\da

sys.modules에 utils 있나?

\da

있음

\da

기존 module 객체 사용
\end{dhksmry}

그래서 일반적으로 \zsh{utils.py}를 처음부터 다시 실행하지 않는다.

\subsubsection{그래서 같은 module 객체인지 확인할 수도 있다}

\begin{dhkcode}
import sys
import utils

print(utils is sys.modules["utils"])
\end{dhkcode}

결과:

\begin{dhkzsh}
True
\end{dhkzsh}

\begin{dhksmry}[centered]
현재 namespace
utils

\da

module object

\ua

sys.modules["utils"]
\end{dhksmry}

두 이름/참조가 \textbf{동일한 module 객체를 가리키므로}:

\begin{dhkcode}
utils is sys.modules["utils"]
\end{dhkcode}

가 \zsh{True}인 것이다.

\subsubsection{이제 전체 import 과정을 조금 더 정확하게 만들 수 있다}

단순하게:

\begin{dhksmry}[centered]
import

\da

sys.path 탐색

\da

module 발견
\end{dhksmry}

이라고 이전에 배웠던 것이 이제 한 단계 정밀해진다.

\begin{dhksmry}[centered]
import src.utils

\da

이미 로드되어 있는가?

\da

sys.modules 확인

\da

No

\da

import system이 module 탐색

\da

sys.path

\da

module 발견

\da

module 코드 실행
\end{dhksmry}

\subsubsection{그리고 이것이 \zsh{__name__}과도 연결된다}

\zsh{utils.py}가 import 되면서 실행될 때

Python은 module namespace 에 여러 특별한 이름도 준비한다.

개념적으로:

\begin{dhksmry}
    utils module namespace


    __name__ \ra "utils"

    __package__ \ra ...

    __file__ \ra ".../utils.py"

    add \ra function object

    x \ra 10

    ...
\end{dhksmry}

그래서 \zsh{utils.py}내부 코드가 실행되는 동안에도:

\begin{dhkcode}
print(__name__)
\end{dhkcode}

이 가능했던 것이다.

\zsh{__name__}이라는 마법 같은 함수가 따로 있는게 아니라

\textbf{Python이 module namespace에 준비해둔 특별한 이름}이라고 보면 된다.

여기까지 이해하면 다음에는

\begin{dhksmry}
    import utils와 from utils import add가 

    객체/binding 관점에서 정확히 어떻게 다른지

    그리고 왜 import한 변수는 원본 module에서 나중에 바뀌었는데

    내쪽에서는 안바뀌지? 정도까지 들어가면 된다.
\end{dhksmry}

\subsubsection{\zsh{import ... / from ... import ...} $\rightarrow$ 객체와 바인딩 관점에서 분리}

먼저 모듈 하나를 만들자.

\zsh{config.py}:

\begin{dhkcode}
learning_rate = 0.001

def show_lr():
    print(learning_rate)
\end{dhkcode}

다른 파일 \zsh{train.py}에서:

\begin{dhkcode}
import config
\end{dhkcode}

한다고 하자.

앞에서 배운 과정에 따라 \zsh{config.py}가 처음 import 되면

실행되고 module namespace 가 만들어진다.

\begin{dhkdir}
config module object
> namespace
>> learning_rate → 0.001
>> show_lr → function object
\end{dhkdir}

그리고 \zsh{train.py}쪽에서는:

\begin{dhksmry}
    trian.py namespace
    
    config \ra config module object
\end{dhksmry}

즉 \zsh{import config}가 현재 namespace 에 만드는 핵심 이름은 \zsh{config}이다.

그래서:

\begin{dhkcode}
config.learning_rate
config.show_lr()
\end{dhkcode}

라고 접근한다.

\subsubsection{\zsh{from config import learning_rate}는 다르다}

이번에는:

\begin{dhkcode}
from config import learning_rate
\end{dhkcode}

라고 해보자.

\zsh{config module}을 찾고 로드하는 과정 자체는 여전히 필요하다.

하지만 마지막에 현재 namespace 로 가져오는 이름이 다르다.

\begin{dhkbranch}
node main :: config module namespace // learning_rate → 0.001
arrow main :: from config import learning_rate
node main :: train.py namespace // learning_rate → 같은 0.001 객체
\end{dhkbranch}

즉 현재 namespace 에:

\begin{dhkzsh}
config
\end{dhkzsh}

라는 이름을 만드는 게 아니라:

\begin{dhkzsh}
learning_rate
\end{dhkzsh}

라는 이름을 직접 binding 하는 것이다.

그래서:

\begin{dhkcode}
print(learning_rate)
\end{dhkcode}

라고 바로 쓸 수 있다.

\subsubsection{여기서 중요한 실험}

\zsh{config.py}:

\begin{dhkcode}
learning_rate = 0.001
\end{dhkcode}

\zsh{train.py}:

\begin{dhkcode}
import config

print(config.learning_rate)

config.learning_rate = 0.01

print(config.learning_rate)
\end{dhkcode}

결과:

\begin{dhkzsh}
0.001
0.01
\end{dhkzsh}

이건 어렵지 않다.

\begin{dhksmry}[centered]
config

\da

module object

\da

learning_rate \ra 0.001
\end{dhksmry}

에서 module namespace 의 \zsh{learning_rate} binding을:

\begin{dhksmry}[centered]
learning_rate \ra 0.01
\end{dhksmry}

로 바꾼 것이기 때문이다.

\subsubsection{그런데 이렇게 하면 재미있는 일이 생긴다}

\begin{dhkcode}
from config import learning_rate

print(learning_rate)

learning_rate = 0.01

print(learning_rate)
\end{dhkcode}

당연히:

\begin{dhkzsh}
0.001
0.01
\end{dhkzsh}

이다.

그런데 이것은 \textbf{\zsh{config.learning_rate}를 0.01로 바꾼 것이 아니다}.

처음에는:

\begin{dhksmry}[centered]
config namespace
learning_rate

\da

0.001

\ua

train namespace
learning_rate
\end{dhksmry}

두 namespace 의 서로 다른 이름이 같은 객체를 가리키고 있었던 것이다.

그 후 \zsh{train.py}에서:

\begin{dhkcode}
learning_rate = 0.01
\end{dhkcode}

하면 \textbf{train namespace의 binding만 바뀐다}.

\begin{dhksmry}
    config namespace
    learning_rate \ra 0.001

    train namespace
    learning_rate \ra 0.01
\end{dhksmry}

\subsubsection{반대로 \zsh{import config}는 module 을 통해 계속 접근한다}

\begin{dhkcode}
import config

config.learning_rate = 0.01
\end{dhkcode}

에서는 구조가:

\begin{dhkdir}
train namespace

config
| config module object
>learning_rate -> 0.01
\end{dhkdir}

이다.

별도의 \zsh{learning_rate}이름을 \zsh{train} namespace에 만든 게 아니라

\textbf{config module namespace에 들어가서 그 binding을 변경}한 것이다.

그래서 이 둘은 꽤 중요한 차이가 있다.

\begin{dhkcode}
import config
config.learning_rate
\end{dhkcode}

    vs.

\begin{dhkcode}
from config import learning_rate
learning_rate
\end{dhkcode}

첫 번째는 매번 \zsh{config} module의 attribute를 보고 있고,

두 번째는 import 시점에 해당 객체를 \textbf{현재 namespace의 별도 이름에 binding한 것}이다.

\subsubsection{함수도 완전히 똑같다}

\zsh{utils.py}:

\begin{dhkcode}
def add(a,b):
    return a + b
\end{dhkcode}

방법 A

\begin{dhkcode}
import utils

utils.add(10, 20)
\end{dhkcode}

구조:

\begin{dhkdir}
현재 namespace

utils -> @module object
> add -> function object
\end{dhkdir}

방법 B

\begin{dhkcode}
from utils import add

add(10, 20)
\end{dhkcode}

구조:

\begin{dhksmry}[centered]
utils module

\da

add

\da

function object

\ua

add

\ua

현재 namespace
\end{dhksmry}

즉 \textbf{함수 코드를 복사해온 것도 아니다}.

같은 function object를 현재 namespace의 \zsh{add}라는 이름에서도 가리키게 만든 것이다.

\subsubsection{\zsh{as}도 이제 별 것 아니다}

\begin{dhkcode}
import pandas as pd
\end{dhkcode}

이것도 특별한 마법처럼 볼 필요가 없다.

대략:

\begin{dhksmry}[centered]
pandas module object
\ua
pd
\end{dhksmry}

현재 namespace에서 그 module 객체를 \zsh{pd}라는 이름으로 사용하겠다는 것이다.

그래서:

\begin{dhkcode}
pd.read_csv(...)
\end{dhkcode}

가 되는 것이다.

마찬가지로:

\begin{dhkcode}
from pathlib import Path as P
\end{dhkcode}

라면 현재 namespace에:

\begin{dhksmry}[centered]
P \ra Path class object
\end{dhksmry}

가 생기는 셈이다.

\subsubsection{여기까지 하나로 정리}

\begin{dhkdir}
import utils // 현재 namespace
| utils -> @module
> add
> ...
\end{dhkdir}

\begin{dhkdir}
from utils import add // 현재 namespace
| add -> function object
| utils module
> add -> function object(same)
\end{dhkdir}

결국 import 문법도:

\begin{dhksmry}[centered]
객체를 찾고 \ra 어떤 namespace의 어떤 이름에 binding 할 것인가
\end{dhksmry}

라는 원리 위에 있다.

다음에는 여기서 한단계 더 나아가서,

\textbf{mutable 객체를 from ... import ... 했을 때는 왜 조금 다르게 보이는지},

그리고 \textbf{import * 가 왜 좋지 않은지} 살펴본다.

\subsubsection{mutable 객체를 import 했을 때의 동작은?}

이번에는 \zsh{config}에 리스트를 넣어보자.

\zsh{config.py}:

\begin{dhkcode}
hidden_sizes = [512, 256, 128]
\end{dhkcode}

\zsh{train.py}:

\begin{dhkcode}
from config import hidden_sizes

print(hidden_sizes)
\end{dhkcode}

처음 import 직후 구조는:

\begin{dhksmry}[centered]
config namespace

hidden_sizes

\da

[512, 256, 128]

\ua

hidden_sizes

train namespace
\end{dhksmry}

즉 \textbf{리스트를 복사한 것이 아니다}.

두 namespace의 \zsh{hidden_size}라는 이름이 \textbf{동일한 list 객체를 가리킨다}.

따라서:

\begin{dhkcode}
import config
from config import hidden_sizes

print(config.hidden_sizes is hidden_sizes)
\end{dhkcode}

는:

\begin{dhkzsh}
True
\end{dhkzsh}

가 된다.

\subsubsection{여기서 객체 자체를 변경해보자}

\begin{dhkcode}
import config
from config import hidden_sizes

hidden_sizes.append(64)

print(hidden_sizes)
print(config.hidden_sizes)
\end{dhkcode}

결과:

\begin{dhkzsh}
[512, 256, 128, 64]
[512, 256, 128, 64]
\end{dhkzsh}

왜?

\zsh{append()}는 \zsh{hidden_sizes}라는 이름의 binding을 바꾼 것이 아니다.

\textbf{그 이름이 가리키고 있는 list 객체 자체를 변경한 것}이다.

그래서 양쪽에서 변경된 상태가 보인다.

\subsubsection{그런데 재대입하면 완전히 다르다}

이번에는:

\begin{dhkcode}
import config
from config import hidden_sizes

hidden_sizes = [1024, 512]
print(hidden_sizes)
print(config.hidden_sizes)
\end{dhkcode}

결과:

\begin{dhkzsh}
[1024, 512]
[512, 256, 128]
\end{dhkzsh}

왜냐하면:

\begin{dhkcode}
hidden_sizes = [1024, 512]
\end{dhkcode}

는 기존 list를 수정한 게 아니라 \textbf{train namespace의 \zsh{hidden_sizes}라는 이름을 새로운 list 객체에 재바인딩}했기 때문이다.

\begin{dhksmry}
    config namespace
    hidden_sizes \ra [512, 256, 128]


    train namespace
    hidden_sizes \ra [1024, 512]
\end{dhksmry}

이건 mutable/immutable 의 내용과 같은 원리이다.

\subsubsection{따라서 핵심은 import가 아니다}

이게 중요하다.

이 현상을 보고:

\begin{dhksmry}[centered]
\zsh{from import}는 mutable일 때 특별하게 동작한다.
\end{dhksmry}

라고 \textbf{이해하면 안된다}.

특별한 import 규칙이 있는 것이 아니다.

핵심은:

\begin{dhksmry}[centered]
객체 자체 변경(mutaion)

vs.

이름의 재바인딩(rebinding)
\end{dhksmry}

이다.

예를 들어:

\begin{dhkcode}
hidden_sizes.append(64)
\end{dhkcode}

는:

\begin{dhksmry}[centered]
이름

\da

기존 list 객체

\da

객체 내부 상태 변경
\end{dhksmry}

반면:

\begin{dhkcode}
hidden_sizes = [1024, 512]
\end{dhkcode}

는:

\begin{dhksmry}[centered]
    이름 \ra 기존 list (X)

    이름 \ra 새로운 list (O)
\end{dhksmry}

이다.

\subsubsection{dictionary도 똑같다}

\zsh{config.py}:

\begin{dhkcode}
settings = {
    "epochs": 10,
    "batch_size": 32
}
\end{dhkcode}

\zsh{train.py}:

\begin{dhkcode}
import config
from config import settings

settings["epochs"] = 20
print(settings)
print(config.settings)
\end{dhkcode}

둘 다:

\begin{dhkzsh}
{
    "epochs": 20,
    "batch_size": 32
}
\end{dhkzsh}

가 보인다.

왜냐하면:

\begin{dhkcode}
settings["epochs"]=20
\end{dhkcode}

은 dictionary 객체 자체를 변경했기 때문이다.

반대로:

\begin{dhkcode}
settings = {
    "epochs": 100
}
\end{dhkcode}

하면 \zsh{train}쪽 이름만 새로운 dictionary로 재바인딩 된다.

\subsubsection{여기까지 오면 import를 굉장히 단순하게 볼 수 있다}

처음에는:

\begin{dhkcode}
import config
from config imoprt hidden_sizes
\end{dhkcode}

가 Python만의 복잡한 규칙처럼 보이지만, 실제 객체 모델은:

\begin{dhksmry}
    1. module도 객체다.

    2. module은 자신의 namespace를 가진다.

    3. import는 객체를 현재 namespace의 이름에 binding한다.

    4. from ... import ... 도 특정 객체를 
        현재 namespace의 이름에 binding한다.

    5. 그 이후부터는 평번한 Python 객체 규칙이다.
\end{dhksmry}

\begin{dhksmry}[centered]
mutation

\da

같은 객체를 보는 모든 곳에서 관찰됨



rebinding

\da

해당 namespace의 이름만 다른 객체를 가리킴
\end{dhksmry}

이것이 전부이다.

\subsubsection{그리고 \zsh{import *}}

이제 이런 코드도 이해할 수 있다.

\begin{dhkcode}
from config import *
\end{dhkcode}

대략적으로는 \zsh{config}에서 공개 대상으로 취급되는 여러 이름을 

현재 namespace로 가져오는 형태다.

예를 들어:

\begin{dhkcode}
# config.py

learning_rate = 0.001
batch_size =32

def load_config():
    ...
\end{dhkcode}

에서:

\begin{dhkcode}
from config import *
\end{dhkcode}

를 사용하면 현재 namespace에 여러 이름이 한꺼번에 들어올 수 있다.

\begin{dhksmry}[centered]
현재 namespace

learning_rate
batch_size
load_config
...
\end{dhksmry}

문제는 구묘가 커지면:

\begin{dhkcode}
from config import *
from utils import *
from dataset import *
\end{dhkcode}

같은 코드에서

\begin{dhkcode}
load()
\end{dhkcode}

를 봤을 때:

\begin{dhksmry}[centered]
\zsh{load}가 어디서 들어온 것인가?
\end{dhksmry}

를 파악하기 어려워진다는 것이다.

같은 이름이 여러 module에 있으면 namespace 충돌도 생길 수 있다.

그래서 일반적인 애플리케이션 코드에서는 \zsh{import *}를 피하고:

\begin{dhkcode}
import config
\end{dhkcode}

또는:

\begin{dhkcode}
from config import learning_rate, batch_size
\end{dhkcode}

처럼 \textbf{어떤 이름을 가져오는지 명시적으로 표현하는 편이 좋다}.

여기까지 왔으면 \textbf{import/module/package} 학습의 핵심 부분은 거의 끝났다.

지금까지 연결한 범위는:

\begin{dhksmry}[centered]
filesystem

\da

Path / CWD

\da

Python 실행

\da

sys.executable

\da

sys.path

\da

module 탐색

\da

module / package

\da

python file.py vs ptyhon -m package.module

\da

__name__ / __package__

\da

module namespace

\da

sys.modules / import cache

\da

import binding

\da

mutable / mutation / rebinding
\end{dhksmry}
